\chapter{群}

\section{群的定义与基本性质}

\subsection{二元运算}

\begin{definition}[二元运算]
设 $S$ 是一个集合。$S$ 上的一个\emph{二元运算}是映射
\[
  f\colon S\times S\longrightarrow S.
\]
对 $a,b\in S$，通常把 $f(a,b)$ 简记为 $ab$ 或 $a\cdot b$。
\end{definition}

\begin{example}
在 $S=\mathbb R$ 上，映射 $f(a,b)=a+b$ 定义了加法运算。类似地，实数乘法也给出一个二元运算。
\end{example}

\begin{definition}[结合律]
若 $S$ 上的二元运算满足
\[
  f\bigl(a,f(b,c)\bigr)=f\bigl(f(a,b),c\bigr),
  \qquad \forall a,b,c\in S,
\]
则称该运算满足\emph{结合律}。使用乘法记号时，这一等式写作
\[
  a(bc)=(ab)c.
\]
\end{definition}

\begin{theorem}[结合运算下乘积的唯一性]
若 $S$ 上的二元运算满足结合律，则对任意 $n\geq 1$ 及 $a_1,\dots,a_n\in S$，存在唯一确定的乘积
$[a_1,\dots,a_n]$，满足
\begin{enumerate}
  \item $[a_1]=a_1$；
  \item $[a_1,a_2]=a_1a_2$；
  \item 对任意 $1\leq i<n$，
  \[
    [a_1,\dots,a_n]
    =[a_1,\dots,a_i][a_{i+1},\dots,a_n].
  \]
\end{enumerate}
因此，在满足结合律时，多项乘积可以省略括号。
\end{theorem}

\begin{proof}
对 $n$ 作归纳。$n=1,2$ 时结论由定义成立。设长度不超过 $n-1$ 的乘积已经唯一确定，定义
\[
  [a_1,\dots,a_n]=[a_1,\dots,a_{n-1}]a_n.
\]
当 $i=n-1$ 时，所需等式就是上述定义；当 $i<n-1$ 时，由归纳假设和结合律，
\begin{align*}
  [a_1,\dots,a_n]
  &=\bigl([a_1,\dots,a_i][a_{i+1},\dots,a_{n-1}]\bigr)a_n\\
  &=[a_1,\dots,a_i]
    \bigl([a_{i+1},\dots,a_{n-1}]a_n\bigr)\\
  &=[a_1,\dots,a_i][a_{i+1},\dots,a_n].
\end{align*}
而取 $i=n-1$ 又说明任何满足条件的乘积都必须等于
$[a_1,\dots,a_{n-1}]a_n$，故唯一性成立。
\end{proof}

\begin{definition}[交换律]
若
\[
  f(a,b)=f(b,a),\qquad \forall a,b\in S,
\]
则称该二元运算满足\emph{交换律}。使用乘法记号时，这一条件写作 $ab=ba$。
\end{definition}

\subsection{单位元与基本代数结构}

\begin{definition}[单位元]
若存在 $e\in S$，使得
\[
  ea=ae=a,\qquad \forall a\in S,
\]
则称 $e$ 为该二元运算的\emph{单位元}。
\end{definition}

\begin{proposition}
一个二元运算至多有一个单位元。
\end{proposition}

\begin{proof}
若 $e$ 与 $e'$ 都是单位元，则
\[
  e=ee'=e'.
\]
\end{proof}

\begin{definition}[半群与幺半群]
带有一个满足结合律的二元运算的集合称为\emph{半群}；有单位元的半群称为\emph{幺半群}。
\end{definition}

\subsection{群、逆元与幂}

\begin{definition}[逆元]
设 $S$ 是有单位元 $e$ 的半群。若对 $a\in S$，存在 $b\in S$ 使得
\[
  ab=ba=e,
\]
则称 $a$ 可逆，并称 $b$ 为 $a$ 的逆元。
\end{definition}

\begin{proposition}
可逆元素的逆元唯一。
\end{proposition}

\begin{proof}
若 $b$ 与 $b'$ 都是 $a$ 的逆元，则
\[
  b=be=b(ab')=(ba)b'=eb'=b'.
\]
因此可将 $a$ 的唯一逆元记为 $a^{-1}$。
\end{proof}

\begin{definition}[群]
一个带有二元运算的集合 $G$ 称为\emph{群}，若该运算满足：
\begin{enumerate}
  \item 结合律；
  \item 存在单位元；
  \item 每个元素都可逆。
\end{enumerate}
若还满足 $ab=ba$ 对所有 $a,b\in G$ 成立，则称 $G$ 为\emph{阿贝尔群}（或\emph{交换群}）。
\end{definition}

在不致混淆时，群的单位元常记为 $1$。对 $a\in G$ 及 $n\in\mathbb Z$，定义
\[
  a^n=
  \begin{cases}
    \underbrace{a\cdots a}_{n\text{ 个}}, & n>0,\\
    1, & n=0,\\
    (a^{-n})^{-1}, & n<0.
  \end{cases}
\]
由结合律可得
\[
  a^{m+n}=a^m a^n,\qquad \forall m,n\in\mathbb Z.
\]

\begin{definition}[群与元素的阶]
若群 $G$ 的元素个数有限，则称 $G$ 为\emph{有限群}，其元素个数 $|G|$ 称为 $G$ 的\emph{阶}。

对 $g\in G$，若存在 $n\geq 1$ 使得 $g^n=1$，则满足该等式的最小正整数称为 $g$ 的\emph{阶}，记作 $\ord(g)$；若不存在这样的 $n$，则称 $g$ 的阶为无穷。
\end{definition}

\subsection{群中的基本运算法则}

\begin{theorem}[消去律]
设 $a,b,c\in G$。若 $ab=ac$，则 $b=c$；若 $ba=ca$，则同样有 $b=c$。
\end{theorem}

\begin{proof}
由 $ab=ac$，在等式两边左乘 $a^{-1}$，得到
\[
  a^{-1}ab=a^{-1}ac,
\]
故 $b=c$。右消去律同理。
\end{proof}

\begin{theorem}[乘积的逆]
对任意 $a,b\in G$，有
\[
  (ab)^{-1}=b^{-1}a^{-1}.
\]
\end{theorem}

\begin{proof}
直接计算得
\[
  (ab)(b^{-1}a^{-1})=a(bb^{-1})a^{-1}=1,
\]
且
\[
  (b^{-1}a^{-1})(ab)=b^{-1}(a^{-1}a)b=1.
\]
由逆元的唯一性即得结论。
\end{proof}

\input{chapters/01-groups-subgroups}
\section{对称群与置换群}

\subsection{对称群}

设 $T$ 是一个集合，记所有从 $T$ 到自身的映射组成的集合为
\[
  \operatorname{Map}(T,T)=\{f\mid f\colon T\to T\}.
\]
映射的复合给出一个二元运算
\[
  \operatorname{Map}(T,T)\times\operatorname{Map}(T,T)
  \longrightarrow\operatorname{Map}(T,T),
  \qquad (f,g)\longmapsto f\circ g.
\]
复合运算满足结合律，恒等映射 $\operatorname{id}_T$ 是单位元。但是，一般的映射未必可逆，
所以 $\operatorname{Map}(T,T)$ 通常不是群。

\begin{definition}[对称群]
集合 $T$ 上所有双射组成的集合
\[
  \operatorname{Sym}(T)
  =\{f\in\operatorname{Map}(T,T)\mid f\text{ 是双射}\}
\]
在映射复合下构成一个群，称为 $T$ 上的\emph{对称群}。
\end{definition}

当 $T=\{1,2,\dots,n\}$ 时，把 $\operatorname{Sym}(T)$ 记为 $S_n$，称为 $n$ 个文字的对称群。
$S_n$ 的元素称为\emph{置换}，而 $S_n$ 的任意子群称为一个\emph{置换群}。

一个置换 $\sigma\in S_n$ 可以写成双行形式
\[
  \sigma=
  \begin{pmatrix}
    1&2&\cdots&n\\
    \sigma(1)&\sigma(2)&\cdots&\sigma(n)
  \end{pmatrix}.
\]
因为 $\sigma(1),\dots,\sigma(n)$ 是 $1,\dots,n$ 的一个排列，所以
\[
  |S_n|=n!.
\]

\subsection{轮换与对换}

\begin{definition}[轮换]
设 $i_1,\dots,i_d$ 是 $\{1,\dots,n\}$ 中两两不同的元素。若 $\sigma\in S_n$ 满足
\[
  \sigma(i_1)=i_2,\quad \sigma(i_2)=i_3,\quad\dots,\quad
  \sigma(i_{d-1})=i_d,\quad \sigma(i_d)=i_1,
\]
并固定其余所有元素，则称 $\sigma$ 为一个 $d$-\emph{轮换}，记作
\[
  (i_1,i_2,\dots,i_d).
\]
当 $d=2$ 时，轮换称为\emph{对换}。
\end{definition}

在轮换记号中，未出现的文字表示不动点。若两个轮换中出现的文字没有交集，则称它们\emph{不相交}。

\begin{proposition}
两个不相交的轮换可交换。
\end{proposition}

\begin{proof}
两个轮换分别只改变各自所含的文字。对任意文字，先后施行这两个轮换与交换施行顺序所得的结果相同。
\end{proof}

\begin{proposition}[不相交轮换分解]
每个置换都可以写成若干个两两不相交轮换的乘积；除去长度为 $1$ 的轮换并忽略各因子的次序后，这一分解唯一。
\end{proposition}

\begin{proof}
对每个文字反复施行置换。由于集合有限，最终会回到起点，从而得到一个轮换。不同轨道要么相同，要么不相交；所有轨道给出所需分解。轨道由置换本身唯一确定，因此分解也具有所述唯一性。
\end{proof}

\begin{proposition}[对换分解]
每个置换都可以写成若干个对换的乘积。
\end{proposition}

\begin{proof}
只需对每个轮换使用恒等式
\[
  (i_1,i_2,\dots,i_d)
  =(i_1,i_2)(i_2,i_3)\cdots(i_{d-1},i_d).
\]
再结合不相交轮换分解即可。
\end{proof}

\subsection{置换的奇偶性}

\begin{definition}[逆序数与符号]
对 $\sigma\in S_n$，集合
\[
  \{(i,j)\mid 1\leq i<j\leq n,\ \sigma(i)>\sigma(j)\}
\]
称为 $\sigma$ 的\emph{逆序集}，其元素个数称为 $\sigma$ 的\emph{逆序数}，记作
$\operatorname{inv}(\sigma)$。定义 $\sigma$ 的符号为
\[
  \operatorname{sgn}(\sigma)=(-1)^{\operatorname{inv}(\sigma)}.
\]
若 $\operatorname{sgn}(\sigma)=1$，则称 $\sigma$ 为\emph{偶置换}；若符号为 $-1$，则称其为\emph{奇置换}。
\end{definition}

\begin{lemma}
设 $\sigma\in S_n$，$\tau=(i,j)$ 是一个对换，则
\[
  \operatorname{sgn}(\sigma\tau)=-\operatorname{sgn}(\sigma).
\]
\end{lemma}

\begin{proof}
先设 $\tau=(k,k+1)$ 是相邻对换。右乘 $\tau$ 只交换单行表示中第 $k$ 与第 $k+1$ 个位置的值；
这两个值之间的逆序状态发生改变，而它们与其他位置共同贡献的逆序数之和不变。因此总逆序数的奇偶性改变。

一般地，当 $i<j$ 时，
\[
  (i,j)=(i,i+1)\cdots(j-2,j-1)(j-1,j)
  (j-2,j-1)\cdots(i,i+1),
\]
右侧共有 $2(j-i)-1$ 个相邻对换，是奇数个。连续应用前一结论即得所需等式。
\end{proof}

\begin{corollary}
一个置换为偶置换（奇置换）当且仅当它可以表示为偶数个（奇数个）对换的乘积。因此，对换分解中对换个数的奇偶性与分解方式无关。
\end{corollary}

\begin{corollary}
对任意 $\sigma,\tau\in S_n$，有
\[
  \operatorname{sgn}(\sigma\tau)
  =\operatorname{sgn}(\sigma)\operatorname{sgn}(\tau).
\]
\end{corollary}

\begin{definition}[交错群]
$S_n$ 中所有偶置换构成的子群
\[
  A_n=\{\sigma\in S_n\mid \operatorname{sgn}(\sigma)=1\}
\]
称为 $n$ 个文字的\emph{交错群}。
\end{definition}

\begin{corollary}
当 $n>1$ 时，$S_n$ 中奇置换与偶置换各占一半，因而
\[
  |A_n|=\frac{n!}{2}.
\]
\end{corollary}

\begin{proof}
固定一个对换 $\rho$。映射 $\sigma\mapsto\sigma\rho$ 在偶置换与奇置换之间给出一个双射。
\end{proof}

\subsection{共轭公式}

\begin{proposition}[$S_n$ 中的共轭公式]
若 $\sigma$ 的不相交轮换分解为
\[
  \sigma=(i_1,\dots,i_{d_1})
  (i_{d_1+1},\dots,i_{d_2})\cdots,
\]
则对任意 $\tau\in S_n$，
\[
  \tau\sigma\tau^{-1}
  =\bigl(\tau(i_1),\dots,\tau(i_{d_1})\bigr)
   \bigl(\tau(i_{d_1+1}),\dots,\tau(i_{d_2})\bigr)\cdots.
\]
\end{proposition}

\begin{proof}
共轭保持乘积，因此只需考虑 $\sigma=(i_1,\dots,i_d)$。对每个 $1\leq k<d$，
\[
  \tau(i_k)\xmapsto{\tau^{-1}}i_k
  \xmapsto{\sigma}i_{k+1}
  \xmapsto{\tau}\tau(i_{k+1}),
\]
而 $\tau(i_d)$ 被送到 $\tau(i_1)$；其他文字保持不动。故
\[
  \tau\sigma\tau^{-1}
  =\bigl(\tau(i_1),\dots,\tau(i_d)\bigr).
\]
\end{proof}

\subsection{几何体的对称群}

\begin{definition}[几何体的对称群]
设 $\Gamma$ 是欧氏空间中的一个几何体。所有保持 $\Gamma$ 的等距变换在复合下构成一个群，称为 $\Gamma$ 的\emph{对称群}。
\end{definition}

单位圆的线性对称群为
\[
  \mathrm{O}(2,\mathbb R)
  =\{A\in\mathrm{GL}_2(\mathbb R)\mid A^{\mathsf T}A=I\},
\]
其中的变换包括旋转与轴反射。

正 $n$ 边形的对称群记为 $D_n$，称为\emph{二面体群}。给顶点依次标号
$1,2,\dots,n$ 后，每个几何对称都诱导一个顶点置换，因而可把 $D_n$ 视为 $S_n$ 的子群。
一个对称将相邻顶点对 $\{1,2\}$ 送到某个相邻顶点对，而这两个顶点的像唯一决定整个对称；共有 $n$ 种相邻顶点对与两种次序，所以
\[
  |D_n|=2n.
\]

\section{陪集}

\begin{definition}[左陪集与右陪集]
设 $H\leq G$，$a\in G$。集合
\[
  aH=\{ah\mid h\in H\},
  \qquad
  Ha=\{ha\mid h\in H\}
\]
分别称为 $H$ 在 $G$ 中的一个\emph{左陪集}与\emph{右陪集}。
\end{definition}

\begin{example}
在二面体群 $D_n$ 中，设 $\sigma$ 是旋转 $2\pi/n$ 的变换，$\tau$ 是一个反射，并令
\[
  H=\langle\sigma\rangle
  =\{1,\sigma,\dots,\sigma^{n-1}\}.
\]
则 $H$ 是旋转组成的子群，$\tau\notin H$，且
\[
  \tau H=\{\tau,\tau\sigma,\dots,\tau\sigma^{n-1}\}.
\]
此外，$\tau H\cap H=\varnothing$，并且 $\tau H\cup H=D_n$。
\end{example}

\begin{proposition}
设 $a,b\in G$。映射
\[
  f\colon aH\longrightarrow bH,
  \qquad g\longmapsto ba^{-1}g
\]
是双射。因此，$H$ 的任意两个左陪集具有相同的基数。
\end{proposition}

\begin{proof}
若 $g=ah\in aH$，则 $f(g)=bh\in bH$，所以映射良定义。定义
\[
  \varphi\colon bH\longrightarrow aH,
  \qquad g\longmapsto ab^{-1}g.
\]
直接计算可得 $f\circ\varphi=\operatorname{id}_{bH}$ 且
$\varphi\circ f=\operatorname{id}_{aH}$，故 $f$ 为双射。
\end{proof}

\begin{proposition}[左陪集相等的判别]
对任意 $a,b\in G$，
\[
  aH=bH\quad\Longleftrightarrow\quad a^{-1}b\in H.
\]
\end{proposition}

\begin{proof}
若 $aH=bH$，则 $b\in aH$，故存在 $h\in H$ 使 $b=ah$，于是 $a^{-1}b=h\in H$。

反过来，若 $a^{-1}b=h\in H$，则 $b=ah$，所以 $bH\subseteq aH$。又由
$b^{-1}a=h^{-1}\in H$ 可得 $aH\subseteq bH$，故 $aH=bH$。
\end{proof}

